\section{A companion over ten unrestricted integers}
\label{sec:shared-tap-sun}

The shared-kernel-tap source also composes with the unchanged Sun finalizer
of package~59. This gives a fixed polynomial
\[
 F_{10}(f,g,h,k,l,w,x,y,z,m)\in\mathbb Z[f,g,h,k,l,w,x,y,z,m]
\]
of \textbf{exact total degree $7\,428\,120$}, with a division-free evaluator
of \textbf{525 binary arithmetic operations}, such that the inherited
compiler and arithmetic theorems, together with the new source argument,
prove in ZFC
\[
 \neg\exists(f,g,h,k,l,w,x,y,z,m)\in\mathbb Z^{10}\;(F_{10}=0)
 \quad\Longleftrightarrow\quad\operatorname{Con}(\mathrm{ZFC}).
\]
The domain here is all integers. Relative to package~59's ten-integer
polynomial of degree $8\,469\,528$ with 533 gates, the degree decreases by
$12.2959\%$ and eight arithmetic operations are saved. The comparison is
at commit \texttt{beb0192cb8dc36f78c369c080f2adfbc1887a486}.

\subsection{The revised source polynomials}

Use the primary shared-tap certificate, with its complete old31 coefficient
tuple, the original unshifted 31-coordinate positive source, and the two
invertible row operations. The 89 weighted source records use 84 unique
kernel taps. Each shared coefficient is inserted once. The target identities
are
\[
 [Z^{q_j}]\{D_2(Z)U_p(Z)^2+D_1(Z)U_p(Z)+D_0(Z)\}
       =2\rho_j(p),\qquad U_p(Z)=\sum_i p_iZ^{w_i}.
\]
The certified numerical parameters are
\[
 W=116,\quad Q=1841,\quad H=1873,\qquad
 (\deg D_0,\deg D_1,\deg D_2)=(1384,1739,1641).
\]
The leading coefficient of $D_2$ is $+1$. Every kernel coefficient has
absolute value below $C_0=2^{2134927}$. Put
\[
 \kappa=2134940,\qquad \beta_S=2^\kappa>4096C_0.
\]
The even exponent is part of the source binding. Define
\begin{align*}
 b&=1+3f,& B&=\beta_Sb^2,\\
 \ell&=\sum_iB^{w_i},& U_c&=\ell+g,\\
 A_T&=\sum_jB^{q_j},&N_0&=B^{117},\\
 N_1&=B^{1874},&N&=N_0N_1=B^{1991},\\
 M_0&=N_0-1-(b-1)\ell,&J_c&=2G_c+A_T,\\
 S_c&=g+J_cN_0,&T_c&=M_0+(B-2)A_TN_0,\\
 \mathcal C&=bB^{116}.
\end{align*}
Here the remaining expressions are
\begin{align*}
 G_c&=D_2(B)U_c^2+D_1(B)U_c+D_0(B),\\
 X&=N^2+(S_c+T_c)N+N-1-S_c,\\
 Y&=\beta_S^{15}2^{\kappa-21}B^{2114}b
     =2^{2130\kappa-21}b^{4229}.
\end{align*}
These are polynomial abbreviations, with fixed exponents. Their exact
degrees are $7964,4229,233$ for $X,Y,\mathcal C$, respectively.

Suppose first that $b$ is a square, $f\ne0$, and
$0<g<2\mathcal C$. Then $b\ge4$, and
$B^{116}\le U_c<3bB^{116}$. The kernel estimates give
\[
 D_2(B)>\tfrac12B^{1641},\qquad
 |D_1(B)|<2C_0B^{1739},\qquad |D_0(B)|<2C_0B^{1384}.
\]
After normalization by $B^{1641}U_c^2$, the two possibly negative terms
are at most $2C_0/B^{18}$ and $2C_0/B^{489}$, each less than $1/8$.
Consequently $G_c>\tfrac14B^{1641}U_c^2>0$. The upper bound
\[
 |G_c|<26C_0b^2B^{1873},\qquad 2G_c<N_1/64
\]
and the 32-position gap $H-Q$ prove
\[
 0<J_c<N_1/32,\qquad J_c+(B-2)A_T+2<N_1.
\]
Since $0\le M_0<N_0$ and $0<g<N_0$, this yields the strict range
\begin{equation}
 0\le S_c,\qquad 0\le T_c,\qquad S_c+T_c<N.
 \label{eq:new-sun-range}
\end{equation}
It also gives $X\ge N^2>3b$ and $Y>\max(b,256)$. As $\kappa$ is even,
the source congruence is $X\equiv1-22\equiv0\pmod3$.

When $b$ is a power of two and $S_c\mathbin{\&}T_c=0$, the low mask recovers
$g=\sum_i n_iB^{w_i}$ with $0\le n_i<b$. Thus $U_c$ encodes positive
coordinates $p_i=n_i+1$. Every signed convolution coefficient $h_s$ obeys
\[
 |h_s|<993C_0b^2<B/4.
\]
The target padding in $J_c$ and the mask $B-2$ test $h_s=0$ even when
the incoming carry is $-1$. Therefore the same borrow argument proves
\[
 S_c\mathbin{\&}T_c=0
 \quad\Longleftrightarrow\quad
 g\text{ encodes }n_i\in[0,b),\quad \rho_j(n+1)=0\text{ for all }j.
\]
For $b=2^k$ and $B=2^v$, the actual mask population and block exponent are
\[
 s_2(T_c)=139v-31k-22,\qquad \log_2N=1991v.
\]
Hence
\[
 1+\log_2N+s_2(T_c)=2130v-31k-21=\log_2Y.
\]
Using \eqref{eq:new-sun-range}, the exact two-block population identity
and Legendre's formula turn $Y\mid\binom{2X}{X}$ into precisely the source
zero condition. Conversely a positive source zero is not the all-one
vector, since its fifth row has value one there. Choosing arbitrarily
large $b=4^s$ and encoding $n_i=p_i-1$ gives $B\le g<\mathcal C$ and
the required forward source witnesses.

\subsection{The explicit unchanged finalizer}

Here are the fixed polynomial formulas, with local notation distinguished
from the source base and masks:
\begin{gather*}
 L=lY,\quad U=2LX,\quad V=4gwY,\quad W_*=bw,\\
 A=U(V+1),\quad B_*=2X+1,\quad C=B_*+(A-2)h,\\
 P=U^2V,\quad K=X+1+k(P-2),\\
 S_*=2A-5,\quad T_*=3W_*C-2(W_*^2-1),\\
 a_1=(P^2-4)K^2+4,\\
 \Delta_*=8\mathcal C^3gK^2
       -g^2\{32(C-KL)^2\mathcal C^3+g^2K^2\},\\
 O=f^2l^2\Delta_*,\quad \xi=zx,\quad
 a_0=b\{1+(3O-4)b^2z^2\},\\
 D=(A^2-4)C^2+4,\quad E=C^2D\xi,\quad
 F=4(A^2-4)E^2+1,\\
 G_*=2(A+2)(A-2)^2E^2-1,\quad \Omega=16E,\\
 H_*=C+(B_*-C)F+\Omega Fy,\quad
 I_*=(G_*^2-1)H_*^2+1,\quad a_2=DFI_*.
\end{gather*}
Finally set
\begin{align*}
 r&=T_*+S_*y,&\Sigma&=\xi S_*^5,\\
 \Theta&=(r^3-rS_*^2-S_*^3)(r^2+23S_*^2),\\
 w_0&=1,&w_1&=a_0,&w_2&=a_0(P^2+K^2+3),\\
 Z&=\Sigma m-\Theta,&A_i&=\Sigma^2w_i^2a_i\quad(0\le i\le2),\\
 V_n&=Z^2+A_2-A_0-A_1,&
 W_n&=V_n^2+4Z^2A_2-4A_0A_1.
\end{align*}
The promised single polynomial is
\begin{equation}
 \boxed{F_{10}=W_n^2-16V_n^2Z^2A_2.}
 \label{eq:new-shared-sun10}
\end{equation}

The inherited norm, separator, and intersective-quintic argument first
gives $\Sigma\ne0$, all three $a_i$ squares, and $\Sigma\mid\Theta$.
In particular $z\ne0$. These conclusions use no source positivity.
The coprime factors of the square $a_0$ must be signed squares. Since
$b\equiv1\pmod3$, it cannot be a negative square, so $b$ is a positive
square. The case $b=1$ forces $f=O=0$ and $a_0=1-4z^2<0$, impossible.
Thus $b\ge4$. The second factor being a square forces $O>0$, since for
$O\le0$ it is at most $1-4b^2z^2<0$.

Now $O>0$ implies $\Delta_*>0$. As $\mathcal C>0$, its displayed formula
rules out $g\le0$ and $K=0$; division by $gK^2>0$ gives
$g^3<8\mathcal C^3$. Hence $0<g<2\mathcal C$ before the source estimates
are invoked. This establishes the correct order of the signed bootstrap
and explains why no additional positive padding term is needed.

Sun's inherited arithmetic bridge then recovers $b$ a power of two and
the binomial divisibility, giving a positive source zero. The converse
uses the cofinal forward witnesses, the same Pell bridge, and residue
interpolation. The precise dependencies are Sun's 2021 Theorems
4.12--4.13~\cite{sun2021} and their specialization in his 2026
Propositions 6.1--6.2~\cite{sun},
as preserved in the package59 proof. Thus the new part is the verified
source replacement; the compiler and finalizer theorem are inherited.

\subsection{Exact degree certificate and finite checks}

The upper bound is certified by induction at every gate of the literal
division-free circuit: leaves have degrees zero or one, a product uses
the sum of the two bounds, and a sum or difference uses their maximum.
For this circuit the final bound is
\[
 456(7964)+896(4229)+24(233)+1760=7\,428\,120.
\]
An analytic leading-coefficient certificate gives the matching lower
bound. Under the diagonal specialization, the source leading coefficients
are
\begin{align*}
 x_0&=3^{7964}\beta_S^{3982},&
 y_0&=3^{4229}\beta_S^{2129}2^{\kappa-21},&
 c_0&=3^{233}\beta_S^{116}.
\end{align*}
The scaled radicand $A_2$ has degree $1\,857\,030$ and leading coefficient
$2^{362}3^8x_0^{114}y_0^{224}c_0^6$. It uniquely dominates $V_n$, while
$A_2^2$ dominates $W_n$. The final subtracted term has degree
$5\,735\,370<7\,428\,120$, so $A_2^4$ is the leading term of $F_{10}$.
Its diagonal leading coefficient is
\[
 2^{1448}3^{32}x_0^{456}y_0^{896}c_0^{24}
   =2^{3727056\kappa-17368}3^{7426392}>0.
\]
Thus no cancellation can lower the certified degree.

Independently, after setting all ten variables equal to one indeterminate
$t$, the verifier propagates the coefficient of the certified upper degree
through every gate modulo a prime. It uses the actual fixed coefficients,
authenticated against their complete compressed binary binding. The final
leading coefficients are
\[
 \begin{array}{c|ccc}
 p&101&1009&1000003\\\hline
 [t^{7428120}]F_{10}(t,\ldots,t)\bmod p&95&646&515853.
 \end{array}
\]
These residues also agree with the displayed analytic coefficient. Any
one nonzero residue is a second certificate of the matching lower bound.
This is an exact leading-coefficient computation, not a sample equality
of polynomial values. Separately, exact sparse integer arithmetic verifies the source
coefficient identities, all 59 local circuit identities, the Lucas
factorization, and the eight-factor norm identity. Those identity checks
establish that the circuit computes \eqref{eq:new-shared-sun10}.

The machine-readable circuit and replay files are in the companion
\texttt{sun/} component. Its evaluator has 337 multiplications, 161
additions, and 27 subtractions. The verifier imports no circuit builder
and passes both ordinary Python and \texttt{python -O}. Finite borrow,
population, and source-assignment tests supplement the unbounded proofs;
they do not decide consistency or replace the inherited arithmetic
theorems.
